\section{GitHub Evidence and Stage Transitions}

\subsection{Repository and Tags}

All work is kept in one repository:
\begin{center}\url{https://github.com/SeraphinaXYX/SC2002-Restaurant-Rush}\end{center}
Each stage is marked with a tag on a commit that compiles and runs on its own
(\texttt{javac -d bin src/restaurantrush/*.java}). Checking out a tag gives exactly the
code that existed at the end of that stage, together with the input scripts and recorded
console output for that stage's runs in \texttt{test-runs/}.

\begin{table}[H]
\centering
\begin{tabularx}{\textwidth}{lX}
\toprule
\textbf{Tag} & \textbf{What the snapshot contains} \\
\midrule
\texttt{stage-1} & 20 source files: Regular customers, the three menu items, \texttt{Waiter}, \texttt{Chef}, \texttt{Cashier},
\texttt{Table}, \texttt{Order}, \texttt{WaitingQueue}, \texttt{Restaurant}, \texttt{ActionResult}, the full six-step turn,
abandonment, payment and both endings. Searching the snapshot for ``VIP'', ``Critic'', ``Host'', ``combo'' and
``Happy Hour'' finds nothing (the only hits are the words ``unhappy departures'' in the end-of-turn display). \\
\texttt{stage-2} & 22 source files: everything in \texttt{stage-1} plus \texttt{VIPCustomer} and \texttt{CriticCustomer}.
The arrivals on turns 7 and 13 become a VIP and a Critic, so there are still six arrivals, and the game asks which stage to play. \\
\texttt{stage-3} & 25 source files: everything in \texttt{stage-2} plus \texttt{Host}, \texttt{ComboMeal} and \texttt{HappyHour}.
This is the build the class diagram in Section~\ref{sec:class} describes. \\
\bottomrule
\end{tabularx}
\caption{Contents of each tagged snapshot}
\end{table}

The tables below were produced from \texttt{git diff --numstat} between the tags, so they match what the
compare links show. Changes to \texttt{README.md}, \texttt{TESTING.md} and \texttt{test-runs/} are left out
of the tables; they only document the new stage.

\subsection{Stage 1 $\rightarrow$ Stage 2}

Compare link:
\begin{center}{\small\url{https://github.com/SeraphinaXYX/SC2002-Restaurant-Rush/compare/stage-1...stage-2}}\end{center}

\begin{table}[H]
\centering
\small
\begin{tabularx}{\textwidth}{l l r X}
\toprule
\textbf{Change} & \textbf{File} & \textbf{+/$-$} & \textbf{Why it was needed} \\
\midrule
New & \texttt{VIPCustomer.java} & 38/0 & Overrides \texttt{getSatisfactionLoss()} (5) and \texttt{applyDiscount()} (10\% off). \\
New & \texttt{CriticCustomer.java} & 57/0 & Loss of 12, always picks the most expensive dish, and reacts to service speed in \texttt{onServed()}. \\
Changed & \texttt{GameConfig.java} & 19/3 & Stores the chosen stage; \texttt{createArrival()} returns a VIP on turn 7 and a Critic on turn 13 from Stage~2. \\
Changed & \texttt{Customer.java} & 9/1 & Adds the empty \texttt{onServed(int waitTurns)} hook, called from \texttt{receiveDish()}, so the Critic can react without the controller knowing it is a Critic. \\
Changed & \texttt{Waiter.java} & 9/2 & \texttt{serve()} passes the wait (serve turn minus order turn) to \texttt{receiveDish()} and reports any reaction in its message. \\
Changed & \texttt{GameView.java} & 23/1 & New \texttt{askStage()} prompt; the banner shows the stage. \\
Changed & \texttt{RestaurantRush.java} & 2/1 & Asks for the stage and passes it to \texttt{GameConfig}. \\
Deleted & none & -- & No files were removed. \\
\bottomrule
\end{tabularx}
\caption{Source files added and changed from Stage 1 to Stage 2 (157 lines added, 8 removed)}
\end{table}

Most of this is the kind of small registration and setup change the brief says to expect: \texttt{GameConfig},
\texttt{GameView} and \texttt{RestaurantRush} only learn that a stage number exists and which customer arrives when.
The one change to an existing rule class is the \texttt{onServed()} hook in \texttt{Customer} and the matching call in
\texttt{Waiter.serve()}. We added it because the Critic's reaction depends on how long the food took, and the
cleanest place to find that out is the moment the dish is served. It is a no-op for Regular and VIP customers,
so their behaviour is unchanged. \texttt{GameController.java} is not in the diff at all: it still calls
\texttt{getSatisfactionLoss()}, \texttt{applyDiscount()} and \texttt{chooseItem()} on whatever \texttt{Customer}
it holds, and the subclass decides what happens.

\subsection{Stage 2 $\rightarrow$ Stage 3}

Compare link:
\begin{center}{\small\url{https://github.com/SeraphinaXYX/SC2002-Restaurant-Rush/compare/stage-2...stage-3}}\end{center}

\begin{table}[H]
\centering
\small
\begin{tabularx}{\textwidth}{l l r X}
\toprule
\textbf{Change} & \textbf{File} & \textbf{+/$-$} & \textbf{Why it was needed} \\
\midrule
New & \texttt{Host.java} & 33/0 & A new \texttt{Staff} subclass whose \texttt{seat()} returns an \texttt{ActionResult}, like the Waiter's. \\
New & \texttt{ComboMeal.java} & 43/0 & Implements \texttt{OrderItem}; two \texttt{MenuItem}s, prep units summed, price 90\% of the sum rounded to the nearest cent. \\
New & \texttt{HappyHour.java} & 53/0 & The one-off pricing policy: used or not, turns remaining, the 20\% discount and the +15 recovery amount. \\
Changed & \texttt{Menu.java} & 31/2 & \texttt{Menu(boolean includeCombos)} builds the six combos; adds \texttt{getCombos()}, \texttt{hasCombos()} and \texttt{getItems()}. \\
Changed & \texttt{Customer.java} & 14/1 & \texttt{pickDefault()} gives a 30\% seeded chance of a combo when the menu has combos; adds \texttt{recoverSatisfaction()}. \\
Changed & \texttt{CriticCustomer.java} & 2/1 & Comment only: explains that a Critic never orders a combo. \\
Changed & \texttt{Waiter.java} & 10/3 & \texttt{takeOrder()} takes the \texttt{HappyHour} (null below Stage~3) and applies it after the customer's own discount, before the price is locked into the \texttt{Order}. \\
Changed & \texttt{GameConfig.java} & 12/0 & \texttt{hasHost()}, \texttt{hasHappyHour()} and \texttt{hasCombos()}, so each stage switches on only its own features. \\
Changed & \texttt{GameController.java} & 83/4 & Creates the Host and Happy Hour for Stage~3, adds \texttt{happyHourStart()} and \texttt{hostTurn()}, and calls \texttt{happyHour.endTurn()} at the end of the turn. \\
Changed & \texttt{GameView.java} & 33/2 & Stage~3 menu line, Happy Hour prompt and status line (turns left), recovery message, combos in the menu list. \\
Changed & \texttt{RestaurantRush.java} & 1/1 & \texttt{askStage(3)} instead of \texttt{askStage(2)}. \\
Deleted & none & -- & No files were removed. \\
\bottomrule
\end{tabularx}
\caption{Source files added and changed from Stage 2 to Stage 3 (315 lines added, 14 removed)}
\end{table}

Again, \texttt{GameConfig}, \texttt{RestaurantRush}, the \texttt{Menu} constructor flag and the new view prompts are
setup changes. The real rule changes are three, and each one is something Stage~3 asks for:
\begin{itemize}[leftmargin=*]
\item \textbf{\texttt{GameController}} gained two new phases. The six steps keep their order; Stage~3 only inserts the
Happy Hour step at the start of the turn and the Host step just before the Waiter. The two new methods sit next to the
existing step methods instead of being mixed into them.
\item \textbf{\texttt{Waiter.takeOrder()}} is where the price is locked, so it is the only correct place to apply Happy Hour.
The customer discount is still applied first, which gives the brief's example of $\$15 \times 0.9 \times 0.8 = \$10.80$.
\item \textbf{\texttt{Customer.pickDefault()}} needed the combo rule. The combo branch is only entered when
\texttt{menu.hasCombos()} is true, so in Stages~1 and~2 the seeded random generator is called exactly as before and the
same seed still gives the same dishes. The replay check below confirms this.
\end{itemize}

\subsection{What Stayed the Same, and How We Checked}

Thirteen of the twenty Stage~1 source files are byte-for-byte identical at \texttt{stage-3}:
\texttt{ActionResult}, \texttt{Cashier}, \texttt{Chef}, \texttt{CustomerState}, \texttt{MenuItem}, \texttt{Order},
\texttt{OrderItem}, \texttt{OrderStatus}, \texttt{RegularCustomer}, \texttt{Restaurant}, \texttt{Staff}, \texttt{Table}
and \texttt{WaitingQueue}. \texttt{VIPCustomer} is also unchanged from Stage~2 to Stage~3. So the order lifecycle,
the Cashier's timing (one payment per turn, oldest first, never on the turn a customer becomes \texttt{READY\_TO\_PAY}),
satisfaction loss, abandonment and the victory check were not edited after Stage~1.

Not editing a file does not prove the behaviour is unchanged, because the controller calls it differently, so we also
replayed the old runs on the final build. We extracted the \texttt{stage-3} snapshot into a clean folder, compiled it,
and fed it each Stage~1 input script with \texttt{--echo}, choosing Stage~1 and seed~42. The scripts at \texttt{stage-3}
are the \texttt{stage-1} scripts with one extra first line (\texttt{1}) to answer the stage menu, which did not exist
in Stage~1. We then compared the output with the transcript recorded at the \texttt{stage-1} tag, ignoring only the
stage menu and the banner line that now names the stage. We did the same for the Stage~2 scripts against the
\texttt{stage-2} recordings.

\begin{table}[H]
\centering
\small
\begin{tabularx}{\textwidth}{l l l X}
\toprule
\textbf{Script} & \textbf{Recorded at} & \textbf{Replayed on} & \textbf{Result} \\
\midrule
\texttt{stage1-victory} & \texttt{stage-1} & \texttt{stage-3} & Identical. Victory: 5 paid, \$63.00, average 74.4. \\
\texttt{stage1-defeat} & \texttt{stage-1} & \texttt{stage-3} & Identical. Defeat: 0 paid, \$0.00, 2 unhappy departures. \\
\texttt{stage1-invalid-input} & \texttt{stage-1} & \texttt{stage-3} & Identical, including every ``explain and ask again'' message. \\
\texttt{stage2-victory} & \texttt{stage-2} & \texttt{stage-3} & Identical, including the VIP's discounted bill and the Critic's reaction. \\
\texttt{stage2-defeat} & \texttt{stage-2} & \texttt{stage-3} & Identical. \\
\bottomrule
\end{tabularx}
\caption{Replaying earlier runs on the final build}
\end{table}

Every turn of these games, including the customers' choices of food, came out the same, so the Stage~3 changes did
not leak into the earlier stages. The exact commands are in \texttt{TESTING.md}.

\subsection{How the Stage 1 Design Took the Later Requirements}

We did not know the exact Stage~2 and Stage~3 rules when we first sketched the classes, but we did know that
customer types and menu items were the parts most likely to grow. Three early decisions paid off.
Making \texttt{Customer} abstract meant VIP and Critic were two new subclasses, and \texttt{GameController} did not
change at all in Stage~2. Having \texttt{Order} depend on the \texttt{OrderItem} interface rather than on
\texttt{MenuItem} meant a \texttt{ComboMeal} could be ordered, cooked and paid for with no change to \texttt{Order},
\texttt{Chef} or \texttt{Cashier}. And because every staff action returns an \texttt{ActionResult}, the \texttt{Host}
reuses the same ``explain why it failed and ask again'' handling as the Waiter instead of needing its own error path.

The place that did grow is \texttt{GameController}, by 83 lines in Stage~3, since it has to schedule the new phases.
We accepted this because the turn order belongs in one place and the new phases are separate methods. We also missed
one thing: the \texttt{onServed()} hook was not in our Stage~1 \texttt{Customer}, so Stage~2 had to add it. If we had
listed ``events a customer might react to'' when designing Stage~1, that change would not have been needed.

% ============================================================
% 5. FUNCTIONAL EVIDENCE
% ============================================================
\newpage
